\part{Métodos de Monte Carlo}

\section{Fundamento: el estimador y su error}
\label{sec:mc-fundamento}

Esta sección deriva los tres resultados que sostienen todo lo que sigue: que el
estimador de Monte Carlo es insesgado, que su error decae como $O(N^{-1/2})$
sin importar la dimensión, y cómo se construye su intervalo de confianza.

\subsection{La integral como valor esperado}

Sea $f$ integrable en $[a,b]$ y sea $I = \int_a^b f(x)\,dx$. Si
$U \sim \mathrm{Unif}(a,b)$, su densidad es $p(x) = 1/(b-a)$ sobre $[a,b]$ y
cero fuera. Entonces, por definición de valor esperado,
\begin{equation}
  \mathbb{E}[f(U)] = \int_a^b f(x)\, p(x)\, dx
  = \frac{1}{b-a}\int_a^b f(x)\,dx = \frac{I}{b-a},
  \label{eq:esperanza}
\end{equation}
de donde
\begin{equation}
  \boxed{\;I = (b-a)\,\mathbb{E}[f(U)]\;}
  \label{eq:integral-esperanza}
\end{equation}

El paso es puramente algebraico, pero cambia la naturaleza del problema:
calcular una integral pasa a ser \emph{estimar una media}. Y una media se
estima promediando muestras.

\subsection{El estimador es insesgado}

Sean $U_1, \dots, U_N$ independientes e idénticamente distribuidas
$\mathrm{Unif}(a,b)$, y defínase
\begin{equation}
  \hat I_N = \frac{b-a}{N}\sum_{i=1}^{N} f(U_i).
  \label{eq:estimador}
\end{equation}

\begin{teorema}[Insesgadez]
\label{teo:insesgado}
$\mathbb{E}[\hat I_N] = I$ para todo $N \ge 1$.
\end{teorema}

\begin{proof}
Por linealidad del valor esperado, que no requiere independencia:
\[
  \mathbb{E}[\hat I_N]
  = \frac{b-a}{N}\, \mathbb{E}\!\left[\sum_{i=1}^{N} f(U_i)\right]
  = \frac{b-a}{N} \sum_{i=1}^{N} \mathbb{E}[f(U_i)].
\]
Como todas las $U_i$ tienen la misma distribución, cada sumando vale lo mismo,
$\mathbb{E}[f(U_i)] = \mathbb{E}[f(U)] = I/(b-a)$ por \eqref{eq:esperanza}.
Hay $N$ sumandos iguales, así que
\[
  \mathbb{E}[\hat I_N] = \frac{b-a}{N} \cdot N \cdot \frac{I}{b-a} = I. \qedhere
\]
\end{proof}

Conviene subrayar qué \emph{no} dice el teorema. Insesgado significa que el
estimador acierta \emph{en promedio sobre infinitas repeticiones}, no que una
corrida particular esté cerca. Cuán cerca está una corrida es justamente lo que
mide la varianza.

\subsection{La varianza y la tasa $O(N^{-1/2})$}

\begin{teorema}[Varianza del estimador]
\label{teo:varianza}
Si $\sigma_f^2 := \mathrm{Var}(f(U)) < \infty$, entonces
\begin{equation}
  \mathrm{Var}(\hat I_N) = \frac{(b-a)^2 \sigma_f^2}{N},
  \qquad\text{y por lo tanto}\qquad
  \mathrm{EE}(\hat I_N) = \frac{(b-a)\,\sigma_f}{\sqrt{N}} .
  \label{eq:varianza}
\end{equation}
\end{teorema}

\begin{proof}
Aquí sí hace falta la independencia. Para variables independientes la varianza
de la suma es la suma de las varianzas, y $\mathrm{Var}(cX) = c^2\mathrm{Var}(X)$:
\[
  \mathrm{Var}(\hat I_N)
  = \left(\frac{b-a}{N}\right)^{\!2} \mathrm{Var}\!\left(\sum_{i=1}^{N} f(U_i)\right)
  = \frac{(b-a)^2}{N^2} \sum_{i=1}^{N} \mathrm{Var}(f(U_i))
  = \frac{(b-a)^2}{N^2} \cdot N\sigma_f^2 ,
\]
que se simplifica a \eqref{eq:varianza}. El error estándar es la raíz.
\end{proof}

El error estándar decae como $N^{-1/2}$. Es una tasa \textbf{lenta}: para
ganar un dígito decimal de precisión hay que multiplicar $N$ por cien. Pero
tiene una propiedad que la vuelve valiosa.

\paragraph{Por qué la tasa no depende de la dimensión.} Obsérvese qué entra en
la fórmula \eqref{eq:varianza}: solamente $N$ y $\sigma_f^2$. La dimensión
\emph{no aparece}. La razón es que la derivación nunca usó que $U$ fuera un
número real: solo usó que las $f(U_i)$ fueran independientes, idénticamente
distribuidas y de varianza finita. Si $U$ es un punto en $\mathbb{R}^d$, cada
$f(U_i)$ sigue siendo un \emph{escalar}, el promedio sigue siendo un promedio
de $N$ números, y el mismo cálculo se aplica palabra por palabra.

El contraste con los métodos determinísticos es lo que da valor a esto. Una
regla de cuadratura sobre una rejilla de $m$ puntos por eje necesita $m^d$
evaluaciones, y su error típico es $O(m^{-k}) = O(N^{-k/d})$ para un método de
orden $k$: la tasa \emph{se degrada} al crecer $d$. Monte Carlo mantiene
$N^{-1/2}$ para todo $d$. Existe una dimensión a partir de la cual Monte Carlo
gana, y la Sección~\ref{sec:mc-alta-dimension} la localiza numéricamente.

Lo que sí puede empeorar con la dimensión es la \emph{constante} $\sigma_f$,
no la tasa. La Sección~\ref{sec:mc-alta-dimension} muestra un caso donde
$\sigma_f$ se vuelve catastrófica.

\subsection{Intervalo de confianza del 95\,\%}

El Teorema del Límite Central establece que, si $\sigma_f^2 < \infty$,
\[
  \frac{\hat I_N - I}{\mathrm{EE}(\hat I_N)} \xrightarrow{\;d\;} \mathcal{N}(0,1)
  \qquad \text{cuando } N \to \infty .
\]
Escribiendo $z_{0.975} \approx 1.96$ para el cuantil que deja $2.5\,\%$ en cada
cola, para $N$ grande
\[
  \mathbb{P}\!\left(-z_{0.975} \le \frac{\hat I_N - I}{\mathrm{EE}} \le z_{0.975}\right)
  \approx 0.95 .
\]
Despejando $I$ dentro de la desigualdad se obtiene el intervalo
\begin{equation}
  \boxed{\;
  \mathrm{IC}_{95\%} = \left[\;
    \hat I_N - 1.96\,\frac{(b-a)S_f}{\sqrt{N}},\;\;
    \hat I_N + 1.96\,\frac{(b-a)S_f}{\sqrt{N}}
  \;\right]}
  \label{eq:ic}
\end{equation}
donde $\sigma_f$, que es desconocida, se reemplaza por su estimador muestral
\begin{equation}
  S_f^2 = \frac{1}{N-1}\sum_{i=1}^N \bigl(f(U_i) - \bar f\bigr)^2 ,
  \qquad \bar f = \frac1N \sum_i f(U_i).
  \label{eq:S}
\end{equation}
El divisor $N-1$ (corrección de Bessel) hace que $S_f^2$ sea insesgado para
$\sigma_f^2$; con divisor $N$ el estimador subestima sistemáticamente la
varianza y el intervalo saldría demasiado angosto.

\paragraph{Qué significa realmente el intervalo.} El $95\,\%$ no es la
probabilidad de que $I$ esté en \emph{este} intervalo: $I$ es una constante, o
está o no está. Es la probabilidad de que el \emph{procedimiento} produzca un
intervalo que atrape a $I$, sobre repeticiones del experimento. Esa lectura es
verificable, y se verifica en la Sección~\ref{sec:mc-integrales}: al partir una
corrida larga en 1000 bloques disjuntos, 946 de los intervalos contienen el
valor verdadero, un $94.6\,\%$ frente al $95\,\%$ nominal.

\paragraph{Dos supuestos que conviene no olvidar.} El intervalo
\eqref{eq:ic} es \emph{asintótico}: descansa en el TLC, así que para $N$
pequeño o integrandos muy asimétricos la cobertura real puede alejarse del
$95\,\%$. Y exige $\sigma_f^2 < \infty$; con integrandos de varianza infinita
el TLC no aplica y el intervalo carece de sentido. Un caso límite aparece en la
Sección~\ref{sec:mc-alta-dimension}, donde en $d = 20$ ningún punto cae dentro
de la región, $S_f = 0$ y el intervalo colapsa a un punto.
